%% kapa: exam=1 round=21 year=2010 period=1 form=A subject=economics subject_ko=경제원론 seq=2 number=42
%% source: Q-Net 감정평가사 기출문제 (cst00309) artlSeq=5024181 file=extracted/1008/5024181/2010년 제21회 감정평가사 1차시험 기출문제/1교시.hwp
%% conversion: kordoc 4.24.51 -> tools/gen1.py (auto); review=proofread-sonnet; flags=none
\begin{question}{42}
대출기간이 동일할 때 은행이 담보대출보다 신용대출에 더 높은 이자율을 적용하는 이유는?
\begin{choices}
  \item 담보대출시 담보권 설정 등의 제비용을 고객이 부담하기 때문이다.
  \item 신용대출을 위한 은행의 자금조달비용이 담보대출의 경우보다 낮기 때문이다.
  \item 정부가 담보대출시보다 신용대출시에 더 높은 세율(稅率)을 부과하기 때문이다.
  \item 신용대출시 은행이 부담하는 채무불이행 위험이 담보대출시보다 크기 때문이다.
  \item 은행이 신용대출에 적용되는 이자율이 하락할 것으로 전망하여 은행 이윤을 안정적으로 관리하기 위해서이다.
\end{choices}
\end{question}
\answer{4}
